\subsection{Finite determination and open dense third-obstruction failure}
\label{sec:ghost-prefix-topology}

\begin{result}{C.1.21}{Proved}{A finite prefix determines every fixed class-two coefficient quotient}
\label{lem:ghost-finite-prefix-determines}
Fix \(\ell\ge5\) prime and \(N\ge1\). Let
\(J_N\) be the least index with
\(v_\ell((J_N+4)!)\ge N\). In an original A.44
prefix containing all indices \(i<J_N\), put
\[
 \begin{split}
 S_{\ell,N}&=\bigl(\{2,3\}\cup
             \{q_i,r_i:i<J_N\}\bigr)\setminus\{\ell\},\\
 D&=N+\max_{p\in S_{\ell,N}}v_\ell(p-1).
 \end{split}                                         \tag{76}
\]
Extend this prefix through an index \(L-1\), with
\(L\ge J_N\) and \(v_\ell((L+4)!)\ge D\).
Every integral class-two section of every ordinary
original-A.44 continuation of that finite prefix
has the same coefficient-\(\ell^N\) class-two
reduction up to geometric conjugacy. After their
abelian cocycles are aligned by the explicit
representatives in the proof, the conjugacy is
central.
\end{result}
\begin{proof}
Let \(a_L,b_L\) be the common prefix products.
For each continuation \(\nu\), choose its
literal integral Kummer cocycles as
\[
 \begin{split}
 a_\nu&=\kappa_\ell(a_L)+
                 \sum_{i\ge L}n_i\kappa_\ell(q_{i,\nu}),\\
 b_\nu&=\kappa_\ell(b_L)+
                 \sum_{i\ge L}n_i\kappa_\ell(r_{i,\nu}),
                 \qquad n_i=(i+4)!.
 \end{split}                                         \tag{77}
\]
The same fixed cocycles are used for \(a_L,b_L\).
All sums converge uniformly, and all tail coefficients
are divisible by \(\ell^D\). Hence for two
continuations
\[
 a_1-a_2\in\ell^D Z^1(G_{\mathbb Q},\mathbb Z_\ell(1)),
 \qquad
 b_1-b_2\in\ell^D Z^1(G_{\mathbb Q},\mathbb Z_\ell(1)).
                                                        \tag{78}
\]
An abelian geometric conjugation, lifted through
the surjective abelianization map of the actual
class-two group, aligns each section with (77).
Its central coordinate is then denoted \(h_\nu\).

At an old positive support prime \(p\in S_{\ell,N}\),
ordinarity gives integral primitives
\(b_\nu=\delta u_\nu\). Their difference satisfies
on a chosen arithmetic Frobenius
\[
 (p-1)(u_1-u_2)
       =(b_1-b_2)(F_p)\in\ell^D\mathbb Z_\ell(1).
                                                        \tag{79}
\]
It follows that
\(u_1-u_2\in\ell^{D-v_\ell(p-1)}\mathbb Z_\ell(1)
\subset\ell^N\mathbb Z_\ell(1)\).
The exact integral tame relation gives
\(h_\nu(T_p)=m_pu_\nu\), with the same ordinary
valuation \(m_p\) in both continuations. Thus
their central inertia values agree modulo
\(\ell^N\). At an old negative support prime
use \(a_\nu=\delta v_\nu\) and
\(h_\nu(T_p)=-m_pv_\nu\); equation (79) is
identical with \(a,v\) in place of \(b,u\).
This includes the fixed primes three and two,
respectively. All divisions in (79) take place
in the integral local Tate module before reduction.

At every support prime outside this finite set,
other than the coefficient prime itself, each
nonzero valuation is divisible by \(\ell^N\).
Its integral central inertia therefore vanishes
modulo \(\ell^N\). At a zero-valuation prime
it is already zero integrally. These statements
apply even to primes present in just one of the
two continuations. Wild and non-\(\ell\) tame
inertia have zero image in the coefficient-\(\ell\)
geometric group. Consequently
\[
 w=\bar h_1-\bar h_2
       \in Z^1(G_{\mathbb Q},\mathbb Z/\ell^N(2))
\]
has zero inertia at every prime different from
\(\ell\). It is a cocycle because the two
abelian cocycle pairs agree literally modulo
\(\ell^N\), by (78). The coefficient action
is unramified away from \(\ell\); hence the
closed normal closure of those inertia groups
also lies in the kernel of \(w\), and \(w\)
factors through \(G_{\mathbb Q,\{\ell,\infty\}}\).

The same coefficient exact sequence used for (68),
now with multiplication by \(\ell^N\), proves
\[
 H^1(G_{\mathbb Q,\{\ell,\infty\}},
                       \mathbb Z/\ell^N(2))=0.
\]
Its adjacent integral groups are zero: first
cohomology injects into the computed absolute
weight-two group, and second cohomology is
\(K_2(\mathbb Z)\{\ell\}=0\).
Thus \(w\) is an actual finite coboundary.
This is precisely the central geometric conjugacy
claimed in the lemma.
\end{proof}

\begin{result}{C.2.9}{Proved}{Open dense failure and finite prefixes forcing exclusion}
\label{cor:ghost-open-dense-class3}
Let \(\mathscr U_3\subset\mathscr A\) consist of
the original A.44 outputs having a nonzero local
third-obstruction class at some coefficient prime
\(\ell\ge5\), detected at some finite coefficient
exponent. Then \(\mathscr U_3\) is open and dense
in \(\mathscr A\), and
\(\mathscr U_3\cap\mathscr A_c\) is open and
dense in \(\mathscr A_c\).
Every member has first unrestricted integral
failure at class three. In particular, every finite
legal prefix has a finite extension whose every
ordinary legal continuation is excluded at class
three. In the compact-passing domain such forcing
prefixes can be chosen within that domain, retaining
the affine abelian cover section and excluding its
metabelian lifts.
\end{result}
\begin{proof}
At \(\ell\ge5\) the global integral class-two
class is unique up to geometric conjugacy, by
Corollary~\ref{cor:ghost-class2-fibre24}. Thus a
nonzero obstruction in that component excludes
every full integral class-two section.

We first check that finite-coefficient comparison
really preserves the obstruction in the direction
used here. For \(\ell\ge5\), reduction of the
actual integral Hall coordinates through class
three modulo \(\ell^N\) is a group quotient:
the multiplication, powering, and arithmetic-action
polynomials have only denominators dividing powers
of \(24\), which are units at \(\ell\).
The class-two quotient is the Heisenberg group
over \(\mathbb Z/\ell^N\); the kernel of the
class-three-to-class-two map is
\((\mathbb Z/\ell^N(3))^2\).
The zero-last-coordinate Hall set lift commutes
with reduction, so its actual two-cocycle defect
is the reduction of the integral third obstruction.
A geometric conjugator in class two lifts to class
three; inner conjugation acts trivially on this
central kernel. Conjugating the lifted set section
therefore changes its defect by a coboundary.
Thus geometrically conjugate class-two reductions
have the same finite third-obstruction class,
globally and after restriction to a local group.

Take an output in \(\mathscr U_3\), and a finite
coefficient exponent \(\ell^N\) detecting its
stated local obstruction. Apply
Lemma~\ref{lem:ghost-finite-prefix-determines},
enlarging the prefix to include that local prime
when it is a fresh support prime. Every continuation
has a geometrically conjugate class-two reduction,
and hence the same nonzero finite local obstruction.
The entire resulting prefix cylinder lies in
\(\mathscr U_3\). This proves openness.

Theorem~\ref{thm:ghost-dense-class3} constructs,
after every finite legal prefix, a continuation
whose local obstruction at coefficient \(\ell\ge5\)
is already nonzero modulo \(\ell^2\). It lies
in \(\mathscr U_3\), proving density. Applying
the openness argument to that continuation gives
the asserted finite forcing extension. The same
construction inside \(\mathscr A_c\), together
with Proposition~\ref{prop:ghost-all-passing-genjac}
and the exact morphism (63), proves all the stated
assertions in the compact-passing domain.
\end{proof}

The remaining outputs are contained in the closed
set \(\mathscr A\setminus\mathscr U_3\), which
has empty interior; its intersection with
\(\mathscr A_c\) also has empty relative interior.
No nonemptiness or emptiness assertion for that
closed set is made. Local finite vanishing there
still leaves the integral quotient-convergence
condition of Proposition~\ref{prop:ghost-class3-residue},
and all higher lifting conditions, to be resolved.
